\documentclass[12pt]{extarticle}    % change for bigger fonts
\usepackage[utf8]{inputenc}         % must include
\usepackage{exam}                   % the exam.sty styles

\title{Exam Template}       % name of the exam
\author{@draketeach}        % author(s)
\date{September 2021}       % date last modified

\begin{document}            % begin the document
\begin{coverpage}{}
    NSW Edumacation Something Something     % school name
    \vspace{2cm}    % vertical space
    
    % exam name and year
    \examnameyear{2021}{Yearly Exam}

    % name of course and size
    \coursename{33pt}{Mathematics Extension 3}
    
    \vspace{1cm}    % vertical space
    
    % general instructions
    \begin{generalinstructions}                        
        \item Reading time -- 5 minutes
        \item Working time -- 25 minutes
        \item Write using black pen
        \item Calculators approved by NESA may be used
        \item A reference sheet is provided at the back of this paper
        \item For questions in Section II, show relevant mathematical reasoning and/or calculations
    \end{generalinstructions}
    
    \vspace{1cm}    % vertical space
    
    % section list and total marks
    \begin{sectionlist}{25}
        % section name, marks, and pages
        \textbf{Section I -- 5 marks} (pages 2--3)   
        \begin{itemize}     % section instructions                  
            \item Attempt Questions 1--5
            \item Allow about 10 minutes for this section
        \end{itemize}
        % section name, marks, and pages
        \textbf{Section II -- 20 marks} (pages 4--5)   
        \begin{itemize}     % section instructions  
            \item Attempt Questions 6--7
            \item Allow about 30 minutes for this section
        \end{itemize}
        % section name, marks, and pages
        \textbf{Section III -- 0 marks} (pages 6--7)   
        \begin{itemize}     % section instructions  
            \item Attempt Questions 6--7
            \item Allow about 30 minutes for this section
        \end{itemize}
\end{sectionlist}

\end{coverpage}       % the cover page is a separate file

% section header
\begin{sectionheader}{Trigonometry}
    \textbf{5 marks\\
    Attempt Questions 1--5\\
    Allow about 10 minutes for this section}
    
    Use the multiple-choice answer sheet for Questions 1--5.
\end{sectionheader}

\begin{multiplechoice}              % multiple choice section

    \item Which diagram best represents the solution set of $x^2-2x-3\ge 0$?                % mc question
    \begin{enumerate}               % mc answers
    \begin{multicols}{2}            % two columns
        \item \mcpic{a.png}{5cm}    % mc answer picture
        \item \mcpic{b.png}{5cm}
        \item \mcpic{c.png}{5cm}
        \item \mcpic{d.png}{5cm}
    \end{multicols}
    \end{enumerate}

    \item Given $f(x)=1+\sqrt{x}$, what are the domain and range of $f^{-1}(x)$?    % mc question
    \begin{enumerate}               % mc answers
        \item $x\ge 0,\ y\ge 0$     % mc answer
        \item $x\ge 0,\ y\ge 1$
        \item $x\ge 1,\ y\ge 0$
        \item $x\ge 1,\ y\ge 1$
    \end{enumerate}    

    \item Which of the following is an anti-derivative of $\displaystyle\frac{1}{4x^2+1}$?
    \begin{enumerate}          
        \item $\displaystyle2\tan^{-1}{\left(\frac{x}{2}\right)}+c$    
        \item $\displaystyle\frac{1}{2}\tan^{-1}{\left(\frac{x}{2}\right)}+c$
        \item $2\tan^{-1}{(2x)}+c$
        \item $\displaystyle\frac{1}{2}\tan^{-1}{(2x)}+c$
    \end{enumerate}    
    
    \pagebreak
    
    \item Maria starts at the origin and walks along all of the vector $2\utilde{i}+3\utilde{j}$ , then walks along all of the vector $3\utilde{i}-2\utilde{j}$ and finally along all of the vector $4\utilde{i}-3\utilde{j}$.

    How far from the origin is she? % mc question
    \begin{enumerate}               % mc answers
        \item $\sqrt{77}$           % mc answer
        \item $\sqrt{85}$
        \item $2\sqrt{13}+\sqrt{5}$
        \item $\sqrt{5}+\sqrt{7}+\sqrt{13}$
    \end{enumerate}  
    
    \item The projection of the vector $\displaystyle\binom{6}{7}$ onto the line $y=2x$ is $\displaystyle\binom{4}{8}$.
    
    The point $(6,7)$ is reflected in the line $y=2x$ to a point $A$. 
    
    What is the position vector of the point $A$? 
    \begin{enumerate}   
        \item $\displaystyle\binom{6}{12}$ 
        \item $\displaystyle\binom{2}{9}$
        \item $\displaystyle\binom{-6}{7}$
        \item $\displaystyle\binom{-2}{1}$
    \end{enumerate}  
    
\end{multiplechoice}                % end of mc section

\pagebreak

% section header
\begin{sectionheader}{Section II}
    \textbf{20 marks\\
    Attempt Questions 6--7\\
    Allow about 30 minutes for this section}
    
    Answer each question in the appropriate writing booklet. Extra writing booklets are available.
    
    In Questions 6--7, your responses should include relevant mathematical reasoning and/or \\calculations.
\end{sectionheader}

\begin{shortanswer}         % short answer section

\begin{bookletquestion}{16} % start booklet question

\begin{enumerate}           % start question parts

    \item Let $P(x)=x^3+3x^2-13x+6$.    % question part
    \begin{enumerate}       % start question part parts
        \item\mrks{1} Show that $P(2)=0$
        \item\mrks{2} Hence, factor the polynomial $P(x)$ as $A(x)B(x)$, where $B(x)$ is a quadratic polynomial.
    \end{enumerate}

    \item\mrks{4} Use the principle of mathematical induction to show that for all integers $n\ge 1$,
    $$1\times 2+2\times 5+3\times 8+...+n(3n-1)=n^2(n+1).$$
    
    \item\mrks{3} The diagram shows the graph of $y=f(x)$.

    \begin{center}  % asymptote diagram
    \begin{asy}
        size(0,7cm);
        draw((-5,0)--(5,0), arrow=Arrow);
        label("$x$", (5,0), align=2S+W);
        draw((0,-4)--(0,5), arrow=Arrow);
        label("$y$", (0,5), align=S+2W);
        label("$O$", (0,0), align=2S+2W);
        draw((-2.5,4)..(-2,0)..(0,-2.5)..(2,0)..(2.5,4));
        label("$-2$", (-2,0), align=SW);
        label("$2$", (2,0), align=SE);
        label("$-3$", (0,-2.5), align=SW);
    \end{asy}
    \end{center}

    Sketch the graph of $\displaystyle y=\frac{1}{f(x)}$.
    
    \questionbreak
    
    \item\mrks{4} By expressing $\sqrt{3}\sin{x}+3\cos{x}$ in the form $A\sin{(x+\alpha)}$, solve\\ $\sqrt{3}\sin{x}+3\cos{x}=\sqrt{3}$, for $0\le x\le 2\pi$.

    \item\mrks{2} Solve $\displaystyle\frac{dy}{dx}=e^{2y}$, finding $x$ as a function of $y$.

    \questionend                % End of question...
    \vspace{2cm}                % vertical space
    \centerbold{Please turn over}   % centered bold text
    
\end{enumerate}                 % end question parts

\end{bookletquestion}           % end booklet question

\pagebreak 

\begin{shortquestion}{4}        % start short question

Consider the triangle shown.
\pic{triangle.png}{7cm}         % insert picture
\begin{enumerate}               % begin question parts
    \item\mrks{2} Find the value of $\theta$, correct to the nearest degree.
    \lines{4}                   % insert lines for writing
    \item\mrks{2} Find the value of $x$, correct to one decimal place.
    \lines{4}                   % insert lines for writing
\end{enumerate}

\end{shortquestion}             % end short question
    
\centerbold{End of paper}
\end{shortanswer}               % end short answer section

\pagebreak
\blanklinedpage                 % blank page with lines
\pagebreak
\blankpage                      % blank page

\end{document}                  % end the document
